% !PRINTED XI
\weberpage{XI}
\inhaltpar{113}{Resolventen sechsten Grades}{426}
\inhaltpar{114}{Ueber die Lösung der Ikosaëdergleichung durch transcendente Functionen}{429}
\inhaltabschnitt{Vierzehnter Abschnitt.}{Gruppen linearer ternärer Substitutionen.}
\inhaltpar{115}{Ternäre lineare Substitutionsgruppe vom 168\textsuperscript{sten} Grade}{433}
\inhaltpar{116}{Pole und Axen der ternären Gruppen}{438}
\inhaltpar{117}{Anwendung auf die Gruppe $G_{168}$. Siebenzählige Pole}{442}
\inhaltpar{118}{Die Hauptaxen}{444}
\inhaltpar{119}{Die drei- und sechszähligen Pole}{448}
\inhaltpar{120}{Die Configuration der Gruppe $G_{168}$}{451}
\inhaltpar{121}{Invariantencurven der Gruppe $G_{168}$}{453}
\inhaltpar{122}{Die erste Invariante der Gruppe $G_{168}$ und die Grundcurve}{454}
\inhaltpar{123}{Die höheren Invarianten}{459}
\inhaltpar{124}{Das volle Invariantensystem}{461}
\inhaltabschnitt{Fünfzehnter Abschnitt.}{Das Formenproblem der Gruppe $G_{168}$ und die Theorie der Gleichungen siebenten Grades.}
\inhaltpar{125}{Die Resolventen des Formenproblems}{466}
\inhaltpar{126}{Reduction der allgemeinen Resolvente siebenten Grades auf die specielle}{471}
\inhaltpar{127}{Permutationsgruppe von sieben Ziffern vom Grade 168}{473}
\inhaltpar{128}{Gleichungen siebenten Grades mit einer Gruppe 168\textsuperscript{sten} Grades}{476}
\inhaltpar{129}{Contragrediente Gruppen}{478}
\inhaltpar{130}{Lösung der Gleichung siebenten Grades mit der Gruppe $P_{168}$ durch das Formenproblem der Gruppe $G_{168}$}{481}
\inhaltpar{131}{Möglichkeit der Bestimmung der Functionen $X_1, X_2, X_3$}{484}
\inhaltbuch{Viertes Buch.}{Algebraische Zahlen.}
\inhaltabschnitt{Sechzehnter Abschnitt.}{Zahlen und Functionale eines algebraischen Körpers.}
\inhaltpar{132}{Definition der algebraischen Zahlen}{489}
\inhaltpar{133}{Ganze algebraische Zahlen}{490}
\inhaltpar{134}{Algebraische Körper}{493}
\inhaltpar{135}{Ganze Functionen in einem algebraischen Körper}{495}
\inhaltpar{136}{Die Functionale eines algebraischen Körpers $\Omega$ und der erweiterte Körper $\overline{\Omega}$}{499}
\inhaltpar{137}{Ganze Functionale}{504}
\inhaltpar{138}{Theilbarkeit. Associirte Functionale. Einheiten}{509}
\inhaltpar{139}{Grösster gemeinschaftlicher Theiler}{512}
